\subsection{Планирование цены и прогнозирование прибыли}

На основе рассчитанных суммарных затрат определяется общая стоимость инновационного продукта $K_{\text{по}}$, которая вычисляется по формуле \ref{eq:k_po}.

\begin{equation}\label{eq:k_po}
  K_{\text{по}} = (\Delta K + K_{\text{вн}})\cdot (1 + D_{\text{приб}}),
\end{equation}
где $\Delta K$ — часть затрат на разработку, приходящаяся на один экземпляр комплекса; \\ $K_{\text{вн}}$ — затраты на внедрение; \\ $D_{\text{приб}}$ — процент прибыли, включённый в цену продукта.  

Частичная стоимость разработки $\Delta K$ вычисляется по формуле \ref{eq:delata_k}.

\begin{equation}\label{eq:delata_k}
  \Delta K = \frac{K}{N}\cdot (1 + H_{\text{ст}}),
\end{equation}
где $K$ — общие затраты на проект; \\ $N$ — количество внедряемых единиц комплекса за год;  \\ $H_{\text{ст}}$ — ставка кредитования. На начало 2026 года ставка по кредиту для бизнеса в России составляет 16\%. 

Планируется выпустить только 1 экземпляр комплекса. Тогда согласно формуле \ref{eq:delata_k}:

$$
  \Delta K = 3017244 \cdot (1 + 0.16) = 3500003 \ \text{руб.}
$$

Так как программное обеспечение создаётся по индивидуальному заказу, процент прибыли $D_{\text{приб}}$ принимается равным 0.2. Тогда цена продукта рассчитывается по формуле \ref{eq:k_po}:

$$
  K_{\text{по}} = 3500003 \cdot (1 + 0.2) = 4200003 \ \text{руб.}
$$

Стоимость продукта полностью соответствует ожиданиям рынка, предположенных в при анализе существующих решений. Соответственно, прибыль можно определить по формуле \ref{eq:c_profit} с учётом НДС, равным 22 \%.

\nopagebreak
\begin{equation}\label{eq:c_profit}
  C_{\text{приб}} = K_{\text{по}} \cdot D_{\text{приб}} \cdot (1 - H_{\text{НДС}}) = 655200 \ \text{руб.}
\end{equation}